\noindent\textbf{Purpose:} Facial-expression annotations can contain disagreement that is lost when training uses a single label. This study examines whether distribution-based supervision improves predictions when a classifier reports only its most confident decisions.
\par\smallskip
\noindent\textbf{Methods:} We compared five Swin-T training targets on FER2013 with FER+ annotations, using three paired seeds and secondary ResNet-18 comparisons. Controls preserved dominant target probabilities or tied maxima. The primary outcome was mean annotation disagreement---the fraction of votes inconsistent with the predicted category---among the 80\% most confident predictions. Additional analyses examined distributional agreement, identical accepted images, and calibration transfer. Evaluation retained categories outside the seven outputs. Frozen models were evaluated on 1,440 RAVDESS clips.
\par\smallskip
\noindent\textbf{Results:} For validation-selected Swin-T, distributional supervision reduced disagreement from 24.06\% to 23.68\% relative to the uniform-minority control; the reduction relative to hard targets was 2.61 percentage points. The corresponding ResNet-18 comparison instead increased disagreement by 0.17 percentage points. Distributional supervision nevertheless achieved lower squared distance than both softening controls in every tested seed, backbone and checkpoint state. Comparisons on identical accepted images depended on which model selected them. External comparisons varied with seed, checkpoint and annotation reference. All 36 nonempty empirical calibration-transfer settings exceeded their risk targets on testing.
\par\smallskip
\noindent\textbf{Conclusion:} In these experiments, improved agreement with crowd-label distributions did not consistently improve confidence-selected category predictions. Selective performance must therefore be interpreted in relation to the accepted images and annotation vocabulary.
